% Prüfe dich – je Verfahrenskette eine Aufgabe, gemischt (zusammenbau v0.9)
\einheitenkopf[pruefe][Prüfe dich]{Prüfe dich}
\setcounter{aufgabe}{43}


% pruefe-dich: terme-e4-k2-s5-v2
\begin{kbaufgabe}{Klammere den größten gemeinsamen Faktor aus.}
\begin{kbblock}
\kbzweispaltig[0.46]{\kbspalte{\kbteil{}{$5a^2 - 10a + 15$}{}\kbantwort{= \kbfeld}}}{\kbraum{2}}
\end{kbblock}
\end{kbaufgabe}


% pruefe-dich: terme-e3-k2-s4-v2
\begin{kbaufgabe}{Löse die Klammer auf.}
\begin{kbblock}
\kbzweispaltig[0.46]{\kbspalte{\kbteil{}{$8x - (3y - 6)$}{}\kbantwort{= \kbfeld}}}{\kbraum{2}}
\end{kbblock}
\end{kbaufgabe}


% pruefe-dich: terme-e2-k5-s4-v2
\begin{kbaufgabe}{Multipliziere.}
\begin{kbblock}
\kbteil{}{$3 \cdot (-6a) =$ \lbfeld}{}
\end{kbblock}
\end{kbaufgabe}


% pruefe-dich: terme-e1-k1-s4-v2
\begin{kbaufgabe}{}
\begin{kbblock}
\kbteil{}{Eine Eintrittskarte kostet für Erwachsene $e$ Euro und für Kinder $k$ Euro. Es kommen 2 Erwachsene und 5 Kinder. Schreibe einen Term dafür, wie viel Euro alle zusammen zahlen.}{}
\kbantwort{\kbfeld}
\end{kbblock}
\end{kbaufgabe}


% pruefe-dich: terme-e2-k4-s6-v2
\begin{kbaufgabe}{Fasse zusammen.}
\begin{kbblock}
\kbteil{}{$3a + a =$ \lbfeld}{}
\end{kbblock}
\end{kbaufgabe}
